\usepackage[fontsize=13]{scrextend}
\usepackage[T1]{fontenc}
\usepackage[utf8]{inputenc}
\usepackage[english]{babel}
\AtBeginDocument{
    \renewcommand{\chaptername}{CHAPTER}
    \renewcommand{\figurename}{Figure}
    \renewcommand{\tablename}{Table}
}
\usepackage[left=3.5cm, right=2cm, top=3.5cm, bottom=3cm]{geometry}
\usepackage{graphicx,float}
\usepackage{array,tabularx}
\usepackage{mathtools}
\usepackage{amssymb,amsmath}
\usepackage{mathptmx}       
\usepackage[varvw]{newtxmath}
\usepackage{multicol}
\usepackage{xcolor}
\usepackage{titlesec,titletoc}
\titleformat{\chapter}[hang]{\bfseries\fontsize{16}{19}\selectfont\raggedright}{\MakeUppercase{\chaptername}\ \thechapter}{1ex}{\MakeUppercase}
\titlespacing*{\chapter}{0pt}{40pt}{20pt}
\titleformat{\section}[hang]{\bfseries\fontsize{14}{17}\selectfont\color{black}}{\thesection}{1.5ex}{}
\titlespacing*{\section}{0pt}{12pt}{6pt}
\titleformat{\subsection}[hang]{\bfseries\itshape\fontsize{14}{17}\selectfont\color{black}}{\thesubsection}{1.5ex}{}
\titlespacing*{\subsection}{0pt}{10pt}{4pt}
\titleformat{\subsubsection}[hang]{\mdseries\fontsize{14}{17}\selectfont\color{black}}{\thesubsubsection}{1.5ex}{}
\titlespacing*{\subsubsection}{0pt}{8pt}{4pt}

\titlecontents{chapter}[0pt]{\bfseries\color{black}}{\MakeUppercase\chaptername\ \thecontentslabel\quad}{}{\normalfont\dotfill\bfseries\contentspage}

\titlecontents{section}[2.5ex]{\color{black}}{ \thecontentslabel\quad}{}{\dotfill\contentspage}

\titlecontents{subsection}[5ex]{\color{black}}{ \thecontentslabel\quad}{}{\dotfill\contentspage}

\titlecontents{subsubsection}[7.5ex]{\color{black}}{ \thecontentslabel\quad}{}{\dotfill\contentspage}

\titlecontents{figure}[0pt]{\color{black}}{\figurename\ \thecontentslabel:\quad}{}{\dotfill\contentspage}
\titlecontents{table}[0pt]{\color{black}}{\tablename\ \thecontentslabel:\quad}{}{\dotfill\contentspage}

\usepackage{enumitem}
\newlist{abbrv}{itemize}{1}
\setlist[abbrv,1]{label=,labelwidth=1in,align=parleft,itemsep=0.1\baselineskip,leftmargin=!}

\usepackage[amsmath,thmmarks]{ntheorem}

\theorembodyfont{\normalfont}
\theoremseparator{.}
\newtheorem{definition}{Định nghĩa}[chapter]
\theorembodyfont{\itshape}
\newtheorem{theorem}[definition]{Định lí}
\newtheorem{lemma}[definition]{Bổ đề}
\newtheorem{corrollary}[definition]{Hệ quả}
\newtheorem{Proposition}[definition]{Mệnh đề}

\theorembodyfont{\normalfont}
\newtheorem{example}{Ví dụ}[chapter]

\theoremheaderfont{\itshape}
\theoremsymbol{\ensuremath{\blacksquare}}
\newtheorem*{proof}{Chứng minh}

\usepackage{csquotes}
\usepackage[style=apa,backend=biber]{biblatex}
\DeclareLanguageMapping{english}{english-apa}
\addbibresource{references.bib}
\usepackage{hyperref}
\hypersetup{
	colorlinks=true,
	linkcolor=black,
	citecolor=red,
	urlcolor=black,
	hypertexnames=false
}

\usepackage{indentfirst}
\setcounter{secnumdepth}{3}
\setcounter{tocdepth}{3}
\renewcommand{\baselinestretch}{1.5}
\setlength{\parindent}{1.27cm}
\allowdisplaybreaks
\emergencystretch=3em
\tolerance=2000
\hbadness=10000

\usepackage[labelfont=bf]{caption}
\numberwithin{figure}{chapter}
\numberwithin{table}{chapter}

\usepackage{subcaption}
\usepackage{tikz}
\usetikzlibrary{arrows.meta,positioning,shapes.geometric,calc}
\usepackage{fancyhdr}
\setlength{\headheight}{15.6pt}
\pagestyle{fancyplain}
\fancypagestyle{plain}{
    \fancyhf{}
    \fancyhead[C]{\thepage}
    \renewcommand{\headrulewidth}{0pt}
}